\documentclass[11pt]{article}
\usepackage[margin=1in]{geometry}
\usepackage{amsmath,amssymb,amsthm,mathtools,booktabs}
\usepackage[hidelinks]{hyperref}

\newtheorem{theorem}{Theorem}
\newtheorem{lemma}{Lemma}
\newtheorem{proposition}{Proposition}
\newtheorem{corollary}{Corollary}
\newtheorem{remark}{Remark}

\newcommand{\ip}[2]{\left\langle #1,#2\right\rangle}
\newcommand{\norm}[1]{\left\lVert #1\right\rVert}
\newcommand{\R}{\mathbb R}
\newcommand{\dd}{\mathrm d}
\newcommand{\D}{\mathrm D}

\title{Exact Geodesic Steps versus Euler Steps in the Self-Concordant Level-Set Analysis}
\author{}
\date{}

\begin{document}
\maketitle

\section{Setup}

Let $\phi$ be a strictly convex $C^3$ barrier with
\[
    g(x):=\nabla \phi(x),
    \qquad
    H(x):=\nabla^2\phi(x).
\]
The Hessian metric defines
\[
    \norm{u}_x^2:=u^\top H(x)u,
    \qquad
    \norm{w}_{x,*}^2:=w^\top H(x)^{-1}w.
\]
We write
\[
    V(x,y):=\ip{x-y}{g(x)-g(y)}
\]
for the Jeffreys divergence. Let $d_R$ denote the Riemannian distance induced by $H$.

Throughout this note we use the standard self-concordance normalization
\[
    |\D^3\phi(x)[p,q,q]|
    \le
    2\norm{p}_x\norm{q}_x^2.
\]
Equivalently, the Hessian metric satisfies the exponential stability estimate
\[
    e^{-2d_R(x,z)}H(x)
    \preceq
    H(z)
    \preceq
    e^{2d_R(x,z)}H(x).
\]
We also use the geodesic coercivity inequality
\[
    d_R(x,y)^2\le V(x,y).
\]

For fixed $x,y$, define
\[
    a:=x-y,
    \qquad
    b:=H(x)^{-1}\bigl(g(x)-g(y)\bigr),
\]
so that
\[
    V(x,y)=\ip{a}{b}_x,
    \qquad
    \ip{p}{q}_x:=p^\top H(x)q.
\]
Define also the asymmetry residual
\[
    U(x,y):=a-b.
\]

\section{Projected tangent and descent identity}

Let $L\subseteq \R^n$ be a fixed linear subspace, and let $L^\perp$ denote its Euclidean orthogonal complement. The projected tangent $\dot x$ is defined by
\[
    x+\dot x\in y+L,
    \qquad
    -g(x)-H(x)\dot x\in -g(y)+L^\perp.
\]
Let $P=P_L^x$ be the $H(x)$-orthogonal projection onto $L$, and let $Q=I-P$.

\begin{lemma}[Projection formula]
The projected tangent satisfies
\[
    \dot x=-Pb-Qa.
\]
Consequently
\[
    a+\dot x=PU,
    \qquad
    b+\dot x=-QU.
\]
\end{lemma}

\begin{proof}
The condition $x+\dot x\in y+L$ is equivalent to
\[
    a+\dot x\in L.
\]
The dual condition says
\[
    -g(x)-H\dot x+g(y)\in L^\perp,
\]
or equivalently
\[
    H(b+\dot x)\in L^\perp.
\]
Since $L^\perp$ is Euclidean orthogonal, this means
\[
    b+\dot x\in L_x^\perp,
\]
where $L_x^\perp$ is the $H(x)$-orthogonal complement of $L$.
Thus $a+\dot x$ is the $L$ component of $a-b=U$, while $b+\dot x$ is minus the $L_x^\perp$ component of $U$. Hence
\[
    a+\dot x=PU,
    \qquad
    b+\dot x=-QU.
\]
Solving either identity gives
\[
    \dot x=PU-a=-Pb-Qa.
\]
\end{proof}

Let $x(t)$ be the exact geodesic with $x(0)=x$ and $\dot x(0)=\dot x$. Let
\[
    v(t):=V(x(t),y),
    \qquad
    E:=\norm{\dot x}_x^2.
\]

\begin{lemma}[Exact descent identity]
The projected tangent gives
\[
    v'(0)=-(E+V(x,y)).
\]
\end{lemma}

\begin{proof}
The Riemannian gradient of $V(\cdot,y)$ at $x$ is
\[
    \nabla_R V(x,y)=a+b.
\]
Therefore
\[
    v'(0)=\ip{a+b}{\dot x}_x.
\]
Using $\dot x=-Pb-Qa$ and orthogonality of $P$ and $Q$,
\[
\begin{aligned}
    -v'(0)
    &=\ip{a+b}{Pb+Qa}_x  \\
    &=\ip{Pa+Pb}{Pb}_x+\ip{Qa+Qb}{Qa}_x \\
    &=\ip{a}{b}_x+\norm{Pb}_x^2+\norm{Qa}_x^2.
\end{aligned}
\]
But
\[
    E=\norm{\dot x}_x^2=\norm{Pb}_x^2+\norm{Qa}_x^2
\]
and $V(x,y)=\ip{a}{b}_x$. Hence $v'(0)=-(E+V(x,y))$.
\end{proof}

\section{Residual and energy bounds}

For standard self-concordance, define
\[
    j(v):=2\bigl(\cosh\sqrt v-1\bigr).
\]

\begin{lemma}[Residual bound]
For all $x,y$,
\[
    \norm{U(x,y)}_x
    \le
    2\bigl(\cosh d_R(x,y)-1\bigr).
\]
Consequently,
\[
    \norm{U(x,y)}_x\le j(V(x,y)).
\]
\end{lemma}

\begin{proof}
Let $R=d_R(x,y)$ and let $\gamma:[0,R]\to\Omega$ be a unit-speed geodesic from $x$ to $y$. Then
\[
    x-y=-\int_0^R\gamma'(s)\,\dd s,
\]
and
\[
    g(x)-g(y)=-\int_0^R H(\gamma(s))\gamma'(s)\,\dd s.
\]
Therefore
\[
    U(x,y)=\int_0^R\left[H(x)^{-1}H(\gamma(s))-I\right]\gamma'(s)\,\dd s.
\]
Set
\[
    A_s:=H(x)^{-1/2}H(\gamma(s))H(x)^{-1/2}.
\]
Exponential stability gives
\[
    e^{-2s}I\preceq A_s\preceq e^{2s}I.
\]
A spectral calculation gives
\[
    \left\|\left[H(x)^{-1}H(\gamma(s))-I\right]\gamma'(s)\right\|_x
    \le
    e^s-e^{-s}.
\]
Integrating,
\[
    \norm{U(x,y)}_x
    \le
    \int_0^R(e^s-e^{-s})\,\dd s
    =2(\cosh R-1).
\]
The $V$-bound follows from $R^2\le V(x,y)$.
\end{proof}

\begin{lemma}[Sharp energy envelope from descent]
Let $v:=V(x,y)$ and $u:=\norm{U(x,y)}_x$. Then
\[
    E
    \le
    B(v,u)
    :=
    \left(
    \frac{\sqrt{4v+u^2}+u}{2}
    \right)^2.
\]
In particular,
\[
    E\le B(v):=
    \left(
    \frac{\sqrt{4v+j(v)^2}+j(v)}{2}
    \right)^2.
\]
\end{lemma}

\begin{proof}
By Cauchy-Schwarz and the descent identity,
\[
    E+v=-v'(0)
    \le
    \sqrt E\,\norm{\nabla_R V(x,y)}_x.
\]
Since $\nabla_R V(x,y)=a+b$ and $U=a-b$,
\[
    \norm{\nabla_R V(x,y)}_x^2
    =\norm{a+b}_x^2
    =4V(x,y)+\norm{U(x,y)}_x^2.
\]
Let $s=\sqrt E$. Then
\[
    s^2+v\le s\sqrt{4v+u^2}.
\]
Solving this quadratic inequality in $s$ gives
\[
    s\le \frac{\sqrt{4v+u^2}+u}{2},
\]
and the result follows. The second bound follows from $u\le j(v)$.
\end{proof}

\begin{remark}[Why this improves the crude energy bound]
The crude projection bound gives $E\le d(x,y)$. Since
\[
    d(x,y)=2V(x,y)+\norm{U(x,y)}_x^2,
\]
it has leading behavior $E\lesssim 2V$ near the target. The sharper descent envelope satisfies
\[
    E\le B(V)=V+O(V^{3/2}),
\]
because $j(V)=V+O(V^2)$. This leading constant improvement is one reason the exact-geodesic basin below is larger than the Euler basin certified by a chord-remainder estimate.
\end{remark}

\section{Exact geodesic work bound}

The second-variation identity for Hessian geodesics is
\[
    v''(t)=2E-\ip{\ddot x(t)}{U(x(t),y)}_{x(t)}.
\]
For standard self-concordance,
\[
    \norm{\ddot x(t)}_{x(t)}\le E.
\]
Thus, whenever $v(t)\le v_0$ on an interval, the residual bound gives
\[
    \norm{U(x(t),y)}_{x(t)}\le j(v_0).
\]

\begin{theorem}[Exact geodesic scalar bound]
Let $v_0=V(x,y)$ and define
\[
    F_{\rm geo}(v):=\frac12 B(v)j(v).
\]
If
\[
    F_{\rm geo}(v_0)\le v_0,
\]
then the exact geodesic step stays in the current sublevel set,
\[
    V(x(t),y)\le v_0,\qquad 0\le t\le1,
\]
and satisfies
\[
    V(x(1),y)\le F_{\rm geo}(v_0).
\]
\end{theorem}

\begin{proof}
For any prefix $s\in[0,1]$, Taylor's formula and the descent identity give
\[
\begin{aligned}
    v(s)
    &=v_0+s v'(0)+\int_0^s(s-t)v''(t)\,\dd t \\
    &\le v_0-s(E+v_0)+Es^2
    +\int_0^s(s-t)\left|\ip{\ddot x(t)}{U(x(t),y)}_{x(t)}\right|\,\dd t.
\end{aligned}
\]
As long as the prefix remains in $V(\cdot,y)\le v_0$,
\[
    \left|\ip{\ddot x(t)}{U(x(t),y)}_{x(t)}\right|
    \le E j(v_0).
\]
Using $E\le B(v_0)$,
\[
    v(s)
    \le
    (1-s)v_0-s(1-s)E+\frac{s^2}{2}B(v_0)j(v_0)
    \le
    (1-s)v_0+s^2 F_{\rm geo}(v_0).
\]
If $F_{\rm geo}(v_0)\le v_0$, the right side is at most $v_0$ for $s\in[0,1]$. A first-exit argument proves the sublevel condition, and at $s=1$ gives $v(1)\le F_{\rm geo}(v_0)$.
\end{proof}

\section{Euler endpoint and the discrete work defect}

Let the explicit Euler endpoint be
\[
    x_E:=x+\dot x.
\]
Define the nonlinear gradient remainder
\[
    R_{\rm nl}:=H(x)^{-1}\bigl(g(x_E)-g(x)-H(x)\dot x\bigr).
\]

\begin{lemma}[Euler work identity]
The Euler endpoint satisfies
\[
    V(x_E,y)=\ip{PU}{R_{\rm nl}}_x.
\]
\end{lemma}

\begin{proof}
From the projection identities,
\[
    x_E-y=a+\dot x=PU.
\]
Moreover,
\[
    H(x)^{-1}(g(x_E)-g(y))
    =b+\dot x+R_{\rm nl}
    =-QU+R_{\rm nl}.
\]
Hence
\[
    V(x_E,y)=\ip{PU}{-QU+R_{\rm nl}}_x.
\]
Since $PU\in L$ and $QU\in L_x^\perp$, the cross term vanishes, giving the result.
\end{proof}

The identity says that $R_{\rm nl}$ is the discrete curvature defect, and that pairing it with $PU$ gives the Euler analogue of the integrated curvature work remaining after the flat cancellation.

We use the standard self-concordant gradient-remainder estimate: if $r=\norm{\dot x}_x=\sqrt E<1$, then
\[
    \norm{R_{\rm nl}}_x
    \le
    \frac{r^2}{1-r}
    =
    \frac{E}{1-\sqrt E}.
\]
Therefore
\[
    V(x_E,y)
    \le
    \norm{U(x,y)}_x\frac{E}{1-\sqrt E}.
\]
Using $\norm{U(x,y)}_x\le j(v)$ and $E\le B(v)$ gives the scalar Euler bound
\[
    F_{\rm Euler}(v)
    :=
    j(v)\frac{B(v)}{1-\sqrt{B(v)}}.
\]
This bound is meaningful when $B(v)<1$.

\section{Certified basins for exact geodesic versus Euler}

The following numerical values are obtained by solving scalar equations involving the functions above.

\begin{center}
\begin{tabular}{lll}
\toprule
Bound & Scalar condition & Largest certified value \\
\midrule
Exact geodesic contraction & $F_{\rm geo}(v)\le v$ & $v\le 0.7633246011$ \\
Exact geodesic normalized quadratic & $F_{\rm geo}(v)\le v^2$ & $v\le 0.4190890308$ \\
Euler Dikin-domain condition & $B(v)<1$ & $v<0.4898377813$ \\
Euler contraction & $F_{\rm Euler}(v)\le v$ & $v\le 0.2323044740$ \\
\bottomrule
\end{tabular}
\end{center}

Thus, under these sufficient estimates, exact geodesic tracing certifies a substantially larger self-concordant attraction region than the Euler endpoint. The comparison is not a statement that the exact geodesic endpoint always has smaller $V$ than every numerical endpoint. Rather, exact geodesic tracing removes the discrete nonlinear-gradient remainder that appears in the Euler work identity, leading to the larger certified basin.

Asymptotically,
\[
    j(v)=v+O(v^2),
    \qquad
    B(v)=v+O(v^{3/2}),
\]
so
\[
    F_{\rm geo}(v)=\frac12 v^2+O(v^{5/2}),
\]
whereas
\[
    F_{\rm Euler}(v)=v^2+O(v^{5/2}).
\]
Both are quadratic-type bounds, but the exact geodesic bound has a better leading constant and a larger certified contraction basin.

\section{Interpretation}

The exact geodesic analysis controls the continuous integrated work term
\[
    -\int_0^1(1-t)\ip{\ddot x(t)}{U(x(t),y)}_{x(t)}\,\dd t.
\]
The Euler analysis controls the discrete work term
\[
    \ip{PU}{R_{\rm nl}}_x.
\]
Both are the nonlinear remainder that survives after the projected tangent cancels the flat contribution through
\[
    v'(0)=-(E+V).
\]

A higher-order or work-certified geodesic integrator can be analyzed by comparing its discrete work defect to the exact geodesic work, or directly by proving a bound of the form
\[
    V(\widehat x,y)
    \le
    \Omega(V(x,y)).
\]
The latter is a discrete Lyapunov proof for the actual numerical map and need not first prove that the numerical endpoint is worse than the exact geodesic endpoint. In general, the difference
\[
    V(\widehat x,y)-V(x_{\rm geo}(1),y)
\]
has no fixed sign.

\end{document}
